\scanpage{3}
\mspage{2}
deux cas envisagés dans b), on voit que leur nombre est 1, donc que $B'$ est local. D'autre part, l'hypothèse $X\quad Y$ plat à fibres séparables est héritée par $\unsure{X'}=\operatorname{Spec}(B')\quad Y'=\operatorname{Spec}(A')$, et les fibres génériques sont isomorphes, ce qui nous permet de remplacer $f:X\quad Y$ par $f':X'\quad Y'$, donc de supposer $Y$ géométriquement unibranche. (Ceci montre d'ailleurs que 19.9.7 est essentiellement contenu dans son corollaire 18.9.11).

Compte tenu de 15.5.6.1. nous sommes ramenés à prouver alors que pour tout $x\quad X$ tel que $f(x)=y$ ne soit pas maximal, $\operatorname{Spec}(\underline{\mathrm O}_{X,x})$ n'est pas disconnecté par $x$. Changeant de notations, on est ramené à prouver que lorsque $\dim Y\quad1$, alors $X$ n'est pas disconnecté par son point fermé $x$. Notons que
\[
\operatorname{prof}B=\operatorname{prof}A+\operatorname{prof}B\otimes_A k
\]
et comme $\operatorname{prof}A\quad1$, puisque $A$ est intègre de dimension 1, on voit que $\operatorname{prof}(B)\quad2$ \uline{sauf} si $\operatorname{prof}B\otimes_A k=0$, ce qui équivaut à $\dim B\otimes_A k=0$ puisque par hypothèse $B\otimes_A k$ est séparable sur $k$, donc \deleted{\unsure{intè}} réduit. Or si $\operatorname{prof}B\quad2$, la conclusion voulue résulte du th. de Hartshorne. \deleted{\unsure{D'autre}}

D'ailleurs, si $A$ est normal, \deleted{\unsure{et de dimension}} et si $\dim A\quad2$, on a $\operatorname{prof}A\quad2$ d'où $\operatorname{prof}B\quad2$, d'où encore la conclusion voulue. Dans le cas où l'anneau $A$ envisagé dans 18.9.7. est universellement \deleted{\unsure{catégnaire|caténaire}, \unsure{(hypothèse}}\inserted{japonais} \inserted{on peut observer que} \deleted{\unsure{envisagée}} dans les réductions précédentes, on peut (sans invoquer 0 23.2.5\unsure{)} se ramener \deleted{\unsure{directe}} tout de suite au cas $A$ normal (en prenant pour $A'$ le normalisé de $A/\underline p$ dans une extension finie donnée de son corps des fractions), et par suite \deleted{on} l'argument qui précède montre qu'on est réduit au cas où $A$ est de plus de dimension 1, donc spectre d'un anneau de val discrète, et où $B\otimes_A k$ est de dimension 0, donc réduit à un corps, puisqu'il est réduit, ce qui implique aussitôt que $B$ est lui-même un anneau de valuation discrète, donc $X-x$ est réduit à un point et par suite connexe.
